\section{Méthodologie}

\subsection{Design de l'étude}

Pour répondre à leurs questions de recherche, Sharevski et Kessell ont choisi de
mener des entretiens qualitatifs avec des questions ouvertes. Le choix est logique
, puisque si l'objectif est de comprendre comment des hacktivistes perçoivent la 
désinformation, autant les laisser parler librement plutôt que de les enfermer 
dans un formulaire à cases à cocher.

Avant de commencer quoi que ce soit, l'étude a reçu l'approbation du comité 
d'éthique de DePaul University. Les entretiens se sont déroulés sur Zoom entre 
2022 et début 2023, chaque session durant environ une heure. Les participants 
pouvaient choisir d'activer leur caméra ou non, et avaient la possibilité de 
passer toute question qui les mettait mal à l'aise. Chaque entretien a été 
enregistré, stocké sur un serveur sécurisé, puis retranscrit manuellement. 
Les participants ont donné leur accord final avant le début de l'analyse.

Un détail méthodologique important mérite d'être souligné. Pour éviter que les 
hacktivistes ne confondent le trolling des débuts d'Internet fait "pour le lulz", 
sans intention malveillante avec les opérations de désinformation organisées 
d'aujourd'hui, les chercheurs ont pris soin d'introduire une définition commune 
de la désinformation au début de chaque entretien. C'est une précaution 
qui fait sens quand on sait que certains participants ont eux-mêmes utilisé ces 
techniques dans leur jeunesse hacktiviste.

Les réponses ont ensuite été analysées par deux chercheurs indépendants, qui ont 
identifié sept thèmes récurrents : les antécédents de la désinformation, les 
représentations mentales qu'en ont les hacktivistes, les pratiques de leaking, 
doxing et deplatforming, les opérations anti-désinformation, les tactiques de 
contre-désinformation, la littératie informationnelle, et enfin le concept de 
misinformation hacktivism. Le niveau d'accord entre les deux codeurs était de 
0,82 sur l'échelle Kappa de Cohen, ce qui témoigne d'une bonne fiabilité de 
l'analyse.

Pour le recrutement, les chercheurs ont utilisé le snowball sampling. Ils ont 
contacté un premier groupe de hacktivistes via leurs réseaux personnels, qui ont 
ensuite recommandé d'autres participants, et ainsi de suite. C'est la méthode qui 
s'impose quand on cherche à atteindre des profils aussi peu accessibles, on ne 
recrute pas des hacktivistes avec une annonce LinkedIn.

\subsection{Population étudiée}

Au final, 22 hacktivistes ont accepté de participer, tous volontaires. Pour être 
inclus dans l'étude, il fallait être majeur, résider aux États-Unis, et avoir été 
actif dans la communauté hacktiviste avant 2014, date à partir de laquelle le 
mouvement a largement déserté les plateformes grand public suite à plusieurs 
affaires judiciaires. Les chercheurs ont vérifié ce critère par leurs propres 
investigations OSINT, s'assurant au passage qu'aucun des participants n'avait 
d'antécédents d'activités criminelles ou nuisibles.

Les profils sont variés. Certains travaillent à exposer l'extrémisme d'extrême 
droite en ligne, d'autres aident à traquer des cybercriminels ou à démonter des 
opérations d'influence étrangères. Quelques-uns font encore du trolling et des 
mèmes "à l'ancienne", d'autres font de la vraie sécurité informatique en analysant 
des ransomwares et en proposant des outils gratuits aux utilisateurs ordinaires. 
La majorité d'entre eux sont dans le jeu depuis longtemps, ayant commencé avec 
les ordinateurs pendant l'enfance ou l'adolescence, souvent par simple curiosité, 
parfois pour se protéger d'harceleurs en ligne.

L'échantillon se compose de 13 hommes (59,1\%), 8 femmes (36,4\%) et 1 personne 
non-binaire (4,5\%). Ce n'est pas un hasard : les chercheurs ont fait un effort 
délibéré pour ne pas produire un échantillon exclusivement masculin, sachant que 
la communauté hacktiviste est historiquement très déséquilibrée sur ce plan. Les 
identités des participants étaient connues des chercheurs, mais toute information 
permettant de les identifier a été supprimée dans la publication, c'était 
d'ailleurs une condition posée par les participants eux-mêmes pour accepter de 
parler.
